% 09_inmunizar.tex
\begin{landscape}

	% =============================================================================
	% CONFIGURACIÓN
	% =============================================================================
	\setlength{\tabcolsep}{2pt}
	\setlength{\extrarowheight}{3pt}
	\fontsize{10pt}{10pt}\selectfont

	% =============================================================================
	% CÁLCULO DE ANCHOS DE COLUMNAS
	% =============================================================================
	\setlength{\availablewidth}{\linewidth}

	% Restamos (6 columnas * 2 lados) = 12 espacios para un ajuste matemático exacto
	\addtolength{\availablewidth}{-12\tabcolsep}

	% Restamos el ancho de las líneas verticales (7 líneas * 0.4pt estándar)
	\addtolength{\availablewidth}{-2.8pt}

	% MACRO HEADER
	\newcommand{\tableheader}{\rowcolor{blue!10}
		\textbf{ESPECIE} & \textbf{Nº DE CASOS} & \textbf{INMUNIZACIÓN} &
		\textbf{FÁRMACO} & \textbf{PRINCIPIO ACTIVO} & \textbf{DOSIS/VÍA DE AD\-MINISTRACIÓN} \\
	}

	\begin{longtable}{
		|C{0.1\availablewidth}    % ESPECIE
		|C{0.1\availablewidth}    % CASOS
		|J{0.26\availablewidth}  % INMUNIZACION
		|L{0.14\availablewidth}  % FARMACO
		|J{0.23\availablewidth}  % PRINCIPIO
		|L{0.17\availablewidth}|   % DOSIS
		% SUMA TOTAL: 1.00
	}

		% ──────────────── ENCABEZADO ────────────────
		\caption{\raggedright Registro de Vacunaciones Aplicadas durante el periodo de internado}
		\label{tab:vacunaciones} \\
		\hline
		\tableheader
		\hline
		\endfirsthead

		% ──────────────── CONTINUACIÓN ────────────────
		\multicolumn{6}{l}{\normalsize\textbf{(Continuación) Cuadro \ref{tab:vacunaciones}}. Registro de Vacunaciones} \\[6pt]
		\hline
		\tableheader
		\hline
		\endhead

		% ──────────────── PIES ────────────────
		\hline
		\multicolumn{6}{r}{\footnotesize\textit{Continúa en la siguiente página \ldots}}
		\endfoot

		%\hline
		\multicolumn{6}{l}{\footnotesize\textit{\textbf{Fuente}: Elaboración propia.}}
		\endlastfoot

		% =============================================================================
		% CONTENIDO
		% =============================================================================

		% CANINO
		\multirow{3}{*}{\textbf{Canino}} & 97 & Parvovirus & Hipradog pv & Virus vivo modificado e inactivado & 1ml/animal SC \\
		\cline{2-6}

		& 36 & Parvovirus, moquillo, hepatitis, laringotraqueitis, traqueobronquitis, parainfluenza, leptospirosis & Hipradog p7 Heptavalente & Virus vivo modificado e inactivado & 1ml/animal SC \\
		\cline{2-6}

		& 1 & Rabia & Rabvac 3TK & Virus inactivado & 1ml/animal SC \\
		\hline

		\textbf{Felino} & 5 & Panleucopenia, rinotraqueitis, calicivirus & Triple felina & Virus inactivado & 1ml/animal SC \\
		\hline

		\textbf{Bovino} & 210 & Rabia & BAGOVAC RABIA & Virus inactivado & 1ml/animal SC \\
		\hline

		\rowcolor{green!10} % Verde claro para el total
		\textbf{TOTAL} & \textbf{349} \\
		\cline{1-2}

	\end{longtable}
	\vspace{-12pt}\normalsize
	% Resumen conciso de la tabla
	El cuadro registra las vacunaciones aplicadas durante el internado. Se documentaron 349 casos: 134 en caninos contra parvovirus, moquillo y rabia; 5 en felinos contra panleucopenia, rinotraqueitis y calicivirus; 210 en bovinos contra rabia. Las vacunas incluyen Hipradog pv (virus vivo modificado), Hipradog p7 (virus vivo modificado), Rabvac 3TK (virus inactivado), Triple felina (virus inactivado), BAGOVAC RABIA (virus inactivado), administradas por vía subcutánea a 1ml por animal.

\end{landscape}
